% Opciones empleadas:
%
%	a4paper -> indica el tamaño del papel, en este caso A4.
%	11pt -> tamaño de la fuente 11 puntos.
%	oneside -> sólo escribimos en una cara del folio.
% 
% Otras opciones interesantes:
%
%	twoside -> escribimos a doble cara.
%	openbib -> para que las referencias bibliográficas tengan un salto de línea entre cada campo de la referencia.
%
\documentclass[a4paper,11pt,twoside]{book}

% configuración de márgenes
\usepackage{setspace}
\doublespacing
\setlength\textwidth{15cm}
\setlength\textheight{22cm}
\setlength\oddsidemargin{1cm}
\setlength\evensidemargin{0cm}

% codificación utf8x y símbolos del idioma español (ñ, acentos, ...)
% el documento será necesario editarlo en utf-8
\usepackage[english]{babel}
\usepackage[utf8x]{inputenc}
\usepackage{morefloats}

% puede que queramos usar el símbolo del euro.
\usepackage{eurosym}

% El paquete fancybox nos permite crear cajas de diferentes estilos con facilidad.
% http://www.ctan.org/get/macros/latex/contrib/fancybox/fancybox.pdf
% http://www.mackichan.com/index.html?techtalk/487.htm~mainFrame
\usepackage{fancybox}
\usepackage{multicol}
\usepackage{amsmath}

% Para incluir subfiguras.
\usepackage{subfigure}

% Para incluir gráficos en JPG => compilar con pdflatex.
%\usepackage[pdftex]{graphicx}

% Para incluir gráficos EPS => compilar con latex.
\usepackage[dvips]{graphicx}
\graphicspath{ {Graphics/} }

% Para escribir en color...
%
% ... cuando compilamos con el comando ``latex''
\usepackage[dvips,usenames]{color}
% ... uando compilamos con el comando ``pdflatex''
% \usepackage[pdftex,usenames,dvipsnames]{color}

% Para realiza hiperenlaces
\usepackage[pagebackref=true]{hyperref}
\hypersetup{
    bookmarks=true,         % barra de marcadores
    unicode=false,          % caracteres non-Latin en marcadores de Acrobat
    pdftoolbar=true,        % mostrar barra de herramientas de Acrobat
    pdfmenubar=true,        % mostrar menú de Acrobat
    pdffitwindow=false,     % ajustar ventana al ancho de página
    pdfstartview={FitH},    % ajustar documento al ancho de página
    pdftitle={Memoria BIBEM},    % título
    pdfauthor={},     % autor
    pdfsubject={},   % tema del documento
    pdfcreator={},   % generador del documento
    pdfproducer={}, % productor del documento
    pdfkeywords={}, % lista de palabras clave
    pdfnewwindow=true,      % enlaces en una nueva ventana
    colorlinks=true,       % false: enlaces en caja; true: enlaces coloreados
    linkcolor=blue,          % color de enlaces internos
    citecolor=green,        % color de enlaces a bibliografía
    filecolor=magenta,      % color de enlaces a ficheros
    urlcolor=blue           % color de enlaces externos
}

% Espaciado y ajuste de márgenes
\usepackage{setspace}
\onehalfspacing
% \doublespacing
\setlength{\textwidth}{15cm}
\setlength{\textheight}{22cm}

% Para evitar lineas huérfanas
\widowpenalty=10000
\clubpenalty=10000

% Paquete fancyhdr -> Para modificar la cabecera y pie de páginas.
% http://tug.ctan.org/tex-archive/macros/latex/contrib/fancyhdr/
\usepackage{fancyhdr}
\pagestyle{fancy}
\fancyhf{}
\fancyhf[HR]{\thepage}
\fancyhf[HL]{\nouppercase\rightmark}

% Package booktabs -> Para mejorar el aspecto de las tablas o cuadros.
% http://www.ctan.org/tex-archive/macros/latex/contrib/booktabs/
\usepackage{booktabs}

% Package rotating -> Para poder girar las tablas y dibujarlas a lo largo
% del folio en vez de a lo ancho.
\usepackage{rotating}

% Packages multicol y multirow, para manejar tablas de filas y columnas múltiples.
\usepackage{multicol}
\usepackage{multirow}

% Para empotrar código con resaltado
\usepackage{listings, color}
%Ahora toca configurar el entorno para el código fuente
\definecolor{gray97}{gray}{.97}
\definecolor{gray75}{gray}{.75}
\definecolor{gray45}{gray}{.45}

\lstset{ frame=Ltb,
     framerule=0pt,
     aboveskip=0.5cm,
     framextopmargin=3pt,
     framexbottommargin=3pt,
     framexleftmargin=0.4cm,
     framesep=0pt,
     rulesep=.4pt,
     backgroundcolor=\color{gray97},
     rulesepcolor=\color{black},
     %
     stringstyle=\ttfamily,
     showstringspaces = false,
     basicstyle=\small\ttfamily,
     commentstyle=\color{gray45},
     keywordstyle=\bfseries,
     %
     numbers=left,
     numbersep=15pt,
     numberstyle=\tiny,
     numberfirstline = false,
     breaklines=true,
   }

% minimizar fragmentado de listados
\lstnewenvironment{listing}[1][]
   {\lstset{#1}\pagebreak[0]}{\pagebreak[0]}

\lstdefinestyle{consola}
   {basicstyle=\scriptsize\bf\ttfamily,
    backgroundcolor=\color{gray75},
   }

\lstdefinestyle{SQL}
   {language=SQL,
   }

\lstdefinestyle{Java}
   {language=Java,
   }

\usepackage{url}


% Personalizamos la separación entre párrafos...
\parskip=6pt

% Personalizamos el identado en la primera línea del nuevo párrafo...
\parindent=10pt

% Establecemos el número máximo de niveles de profundidad en las secciones.
\setcounter{secnumdepth}{3}

% Título
\title{Tittle of the Memory}
% Autor
\author{Petra Franjic \and Gustavo Puig}
% Fecha
\date{\today}


\begin{document}

	% \maketitle sirve para generar automática una portada predefinida, pero para un proyecto fin de carrera
	%	de FIC no sirviría porque no cumple las normas de presentación. Podemos hacer dos cosas:
	% 1. Usarla e ignorar las normas (y asumir las consecuencias que pueda tener)
	% 2. Hacernos una portada en LaTeX que cumpla las normas (menos arriesgado)
	%
        %
% Portada.
%

% Nota: Será más cómodo emplear el comando \maketitle que genera una portada de forma automática, pero 
% no incluye toda la información que es necesario incluir en la memoria de un proyecto de fin de carrera
% de la Facultad de Informática de A Coruña.
%

\graphicspath{ {Graphics/} }

\begin{titlepage}

	\begin{center}

		% Logotipo de la universidad.
		%\includegraphics[width = 6cm] {logocei.pdf}
		\vspace{2cm}

		% Nombre de la facultad, de la universidad y del departamento en que se realiza el PFC.
		{\Large{\textbf{UNIVERSITY OF INFORMATICS OF MADRID - UPM}}}
		\\
		{\it \large{\textbf{ARTIFICIAL INTELLIGENCE DEPARMENT}}}
		\vspace{1cm}

		% Indicamos el nombre de la titulación oficial que hemos cursado con tanto esfuerzo.
		{\large INTELLIGENT SERCHING BASED ON METHAHEURISTICS\\MASTER IN ARTIFICIAL INTELLIGENCE}
		\vspace{1cm}

		% Título
		\textbf{\Large Genetic Algorithms}
		\vspace{7cm}
	\end{center}

	\begin{flushright}
		\begin{tabular}{ll}
			% Nombre del alumno.
			\large{\textbf{Alumni:}}	&
			\large{Petra Franjic, Gustavo Puig} \\

			% Nombre del director/tutor del proyecto.
			\large{\textbf{Professor:}}	&
			\large{Director que corresponda} \\

			% Fecha.
			\large{\textbf{Date:}}	&
			\large{2th November 2015} \\
		\end{tabular}
	\end{flushright}

\end{titlepage}


	% Los proyectos de fin de carrera de FIC han de ir acompañados de una serie de documentos adicionales, algunos
	% 	de ellos obligatorios (certificado, resumen, lista de palabras clave) y otros opcionales (dedicatoria
	%	y agradecimientos).
	%
        %
% Resumen del proyecto de fin de carrera
%

\section*{Abstract:}

The aim of this practical assignment was to implement and analyze a multi-objective genetic algorithm with the goal of obtaining the solution of the magic square problem, which is a combinatorial optimization problem. The algorithm that was chosen for this task is the Non-dominated Sorting Genetic Algorithm II (NSGA-II). Simulations with different parameter values were conducted and the performance of the algorithm was measured. The first part of the documentation explains different aspects of the workings of the algorithm and in the second part the results of the carried out simulations are presented and discussed.
        %
% Palabras clave
%

\section*{Key words List:}

List separated by commas


        %
% Dedicatoria
%
\chapter*{}

\begin{flushright}
	{\it Knight Rider, a shadowy flight into the dangerous world of a man who does not exist. Michael Knight, a young loner on a crusade to champion the cause of the innocent, the helpless in a world of criminals who operate above the law.}
\end{flushright}

        %
% Agradecimientos
%

\chapter*{Gratitudes\footnote{In alphabetic order}}

A D. Profesor que corresponda (Posición del profesor), agradecimiento que se quiera poner.



	% FRONTMATTER: TOC, LOF, LOT y descripción/organización de la memoria.
        \frontmatter

%%% DELETED - INCLUDE IF IT'S NECESSARY
        \tableofcontents
        %\listoffigures
        %\listoftables

        %
% Frontmatter - Introducción. Los miembros del tribunal que juzgan los PFC's tienen muchas más memorias que leer, por lo que
%	agradecerán cualquier detalle que permita facilitarles la vida. En este sentido, realizar una pequeña introducción,
%	comentar la organización y estructura de la memoria y resumir brevemente cada capítulo puede ser una buena práctica
%	que permita al lector centrarse fácilmente en la parte que más le interesa.
%

\chapter[Introduccion]{
	Introduccion
}
Genetic algorithms (GAs) belong to a family of stochastic search methods which is covered by the generic term Evolutionary Computation (EC). The main characteristic of EC techniques is that they computationally simulate the natural evolutionary process. The evolutionary process is modeled by the neo-Darwinian paradigm, which combines Darwinian evolution with Mendelian genetics. Other techniques covered by EC are evolution strategies (ESs) and evolutionary programming (EP).
The EC family of optimization techniques finds its place among other stochastic search and optimization approaches such as Simulated Annealing, Monte Carlo methods and Tabu search. These stochastic approaches were developed as an alternative to solving high-dimensional, discontinuous and multimodal problems, where traditional deterministic search techniques often proved ineffective. The disadvantage of these methods is that they don't guarantee to find the absolute best solution, but they do promise to find an acceptable solution in a reasonable amount of time.
A chromosome or an individual is an encoded candidate solution to a problem. It corresponds to a biological genotype. Chromosomes are typically represented as strings of some predefined alphabet A. The alphabet can be binary or non-binary. Typical examples of non-binary alphabets are real and integer numbers. The phenotype of an individual is the point in the solution space that this individual genotype maps to. A population is a set of individuals. Genes encode a particular element of the candidate solution. The values which a gene can take are called alleles.
***Figure search space solution space***
All evolutionary algorithms share a number of common properties. Firstly, all of them utilize the collective learning process of a population of individuals. Secondly, descendants of individuals are generated by randomized processes intended to model mutation and recombination occurring in nature. Lastly, by means of evaluating individuals in their environment, all individuals are assigned a measure of quality or fitness. This way, it is possible to make a comparison of individual fitness. This is the basis of the process of selection, which favors better individuals to reproduce more often than those relatively worse.
There are several basic differences in the utilization of these principles which make genetic algorithms different from the rest of the techniques of EC. Genetic algorithms emphasize recombination (crossover) as the most important search operator and apply mutation with a very small probability. They also use a probabilistic selection operator and often rely on the binary representation of individuals.
The general principle by which the genetic algorithm works is as follows. There exists a population of individuals where each individual represents one possible solution to the given problem. Each individual can be assigned a fitness measure by calculating the value of the function that is being optimized. The selection operator is used to choose those individuals in the population which are going to become the parents of the next generation. The parents produce children by means of the crossover operator, which emulates genetic recombination. Next the mutation operator is applied on the children. Lastly, the reinsertion operator is used to include children in the population of solutions, which closes the cycle of an algorithm.

Algorithm

General algorithm design
There are two typical genetic algorithm designs, the steady-state genetic algorithm and the generational genetic algorithm.
The steady-state genetic algorithm or incremental algorithm can be thought of as a simpler version of the generational genetic algorithm. At every cycle of the algorithm, two parents are chosen, the crossover operator is applied and a child is formed. The child is subsequently mutated and inserted in the population. Since the population size is maintained constant, the child replaces an individual in the old generation. The individual to be replaced is typically the worst one in the population.
Contrary to the steady-state genetic algorithm, the generational genetic algorithm makes a clear distinction between parents and children. At each algorithm cycle, a new generation, composed exclusively of children is formed and the old generation dies out.
In this practical assignment we have chosen to implement the steady-state algorithm with replacement of the worst individual in the population. The main reason against the generational approach is that absence of elitism would have made the convergence process slower. An additional reason for choosing the steady-state algorithm is the much lower computational cost needed to generate each generation of individuals. The scheme of the algorithm is shown in Figure ? and the pseudocode is shown in figure ?.
*** The scheme of the algorithm ***
*** The pseudocode of the algorithm ***


I'll have to add some modifications to the implementation so that the new generation is formed by choosing best individuals in the set of parents + children


Implementation
Class diagrams + explanations + code


Bibliography:
Redes de neuronas artificiales y computacion evolutiva
EAMOP
handbook of CE
Cupic(home university book)

%
% SECCION
%
%\subsection*{Estructura de la memoria}
%
%\paragraph*{Capítulo 1.}
%Resumen capítulo
%
%\paragraph*{Capítulo 2.}
%Resumen capítulo
%
%\paragraph*{Capítulo 3.}
%Resumen capítulo
%
%\paragraph*{Capítulo 4.}
%Resumen capítulo
%
%\paragraph*{Capítulo 5.}
%Resumen capítulo
%
%\paragraph*{Capítulo 6.}
%Resumen capítulo
%
%\paragraph*{Capítulo 7.}
%Resumen capítulo
%
%\paragraph*{Capítulo 8.}
%Resumen capítulo
%
%\paragraph*{Capítulo 9.}
%Resumen capítulo
%
%\paragraph*{Capítulo 10.}
Resumen capítulo


	% MAINMATTER: El contenido, capítulo a capítulo, de la memoria del PFC.
        \mainmatter

	\chapter[Capitulo 1]{\label{identificadorReferenciaCruzada}
Title of the chapter}

Lorem ipsum dolor sit amet, consectetur adipiscing elit. Integer luctus varius augue vel imperdiet. Maecenas neque turpis, accumsan ut eleifend vel, eleifend gravida est. Duis mattis pulvinar porta. Nunc consequat, felis ut malesuada egestas, ipsum urna sodales enim, eu lobortis mauris tellus in augue. In tristique tortor id nisi placerat et rutrum tortor auctor. Aenean erat elit, faucibus non egestas semper, volutpat sit amet ante. Nam nulla justo, suscipit nec sagittis sed, varius id magna. Duis a lorem id felis eleifend aliquet eget vel arcu. Etiam ultricies diam a magna facilisis ornare. Integer tempor massa quis sem gravida bibendum. Sed massa felis, condimentum sit amet bibendum in, tristique sit amet mauris. Nullam tellus ante, ornare eget mollis ac, gravida id nunc. Donec elit neque, aliquam eget auctor sed, gravida in nibh. Vestibulum sit amet diam a magna tempor bibendum quis a ante. Integer augue sapien, facilisis ac malesuada sed, consectetur id turpis. Quisque a leo dapibus felis aliquet vestibulum vitae vitae urna. Nunc ut felis ut metus gravida ultricies. Cras tincidunt nisi ante. Nunc eu quam interdum dui cursus lacinia. Praesent adipiscing ultrices nulla at posuere.

Aliquam gravida facilisis eleifend. Nunc id lorem nibh. Nulla facilisi. Nunc suscipit lectus quis tellus adipiscing adipiscing. Proin ipsum justo, egestas dictum iaculis et, pretium non diam. Aenean vehicula hendrerit lectus, vel bibendum nisi volutpat quis. Fusce molestie condimentum libero, nec scelerisque nisi mollis et. Aliquam placerat accumsan congue. In scelerisque nulla sit amet mauris consectetur consectetur. Suspendisse mattis mi in sapien egestas sit amet accumsan urna sagittis. Mauris metus massa, porta at rutrum eget, sollicitudin eu risus. Donec non diam erat. Nam sagittis, nibh non fermentum ornare, ipsum est mollis felis, eget vulputate tellus ante at erat. Praesent porta tempus velit at faucibus. Donec rhoncus varius nunc, eget imperdiet ligula blandit eget. Morbi pulvinar, nunc ac porta interdum, mauris ligula sollicitudin sem, et sollicitudin turpis magna eu lectus. Pellentesque nulla dui, feugiat quis pharetra eu, accumsan imperdiet quam. Aenean lacinia porttitor leo vitae hendrerit. Pellentesque dictum ultricies dui, et rhoncus nulla sagittis ut. Suspendisse suscipit, risus sed posuere rhoncus, neque metus scelerisque diam, eu gravida quam orci eu dui.

Nulla sit amet sem massa, aliquet auctor magna. Vivamus quis magna sollicitudin felis bibendum mollis. Etiam elementum lacus id quam tincidunt pulvinar. Ut vehicula aliquam ante, vel venenatis libero bibendum eget. Sed gravida lobortis velit, et ornare nunc ullamcorper et. Morbi neque sapien, mattis id semper in, varius at arcu. Duis aliquam ornare sem, eu sagittis libero facilisis non. Nunc quis ultricies augue. Pellentesque ac nulla libero, ut suscipit dui. Morbi aliquet quam non metus iaculis quis accumsan ipsum consectetur.

Mauris mauris odio, varius quis accumsan bibendum, vulputate vel tellus. Sed non consectetur eros. Aenean et mi ut est tincidunt commodo. Donec nunc metus, facilisis a consequat vitae, dignissim quis leo. Cras condimentum, felis a molestie viverra, massa eros consequat nisl, at porta nibh sapien condimentum ante. Duis convallis velit quis metus pulvinar convallis. Suspendisse sed purus non purus dictum ultricies in vitae nisl. Cras convallis erat eu ipsum molestie fermentum. Mauris turpis nunc, condimentum non aliquam ac, ornare at lectus. Donec scelerisque fermentum malesuada. Praesent urna nisl, iaculis sed pharetra eu, tincidunt et lectus. Nulla posuere eros at elit vehicula laoreet. Sed ut erat orci, ac pellentesque nibh. Nulla at sapien sed justo luctus laoreet.

Sed convallis massa laoreet mi venenatis eget dapibus orci varius. Pellentesque purus nisi, placerat nec pretium ac, pretium vitae augue. Vivamus eleifend vestibulum mi, ac dignissim urna vulputate vel. Integer volutpat neque purus. Nulla hendrerit tempus est nec aliquet. Ut facilisis bibendum nisl eget porttitor. Cum sociis natoque penatibus et magnis dis parturient montes, nascetur ridiculus mus. Fusce consectetur sem eget ante rhoncus nec mollis lorem volutpat. Fusce in libero lectus, ut gravida lectus. Nam nulla nunc, dignissim sed sollicitudin a, commodo vel felis. Donec lacinia justo quis neque tempor interdum ut cursus mi. Phasellus venenatis porttitor mi consectetur egestas. Pellentesque quis elit nibh. Pellentesque molestie turpis ac arcu gravida sed ultricies felis pulvinar.
	%%
% T�TULO DEL CAP�TULO
%
\chapter[Fusce vulputate suscipit ipsum]{
	Fusce vulputate suscipit ipsum
	\label{nombre_referencia_al_capitulo_020}
}

Proin sit amet est. Nunc vel sem sed ligula accumsan cursus. Nunc tincidunt purus non elit accumsan vestibulum. Proin ultrices metus non nisl. Duis feugiat enim et sapien. Suspendisse lectus. Aliquam erat volutpat. Vestibulum ullamcorper. Nullam augue urna, pharetra nec, feugiat at, pellentesque a, nibh. Curabitur nulla.


%
% SECCION - T�tulo de la secci�n
%
\section[Nullam ante lorem placerat et]{
	Nullam ante lorem placerat et\footnote{Ejemplo de pie de p�gina desde un t�tulo de secci�n.}
}


Vestibulum pretium, libero in cursus posuere, metus nisi viverra libero, sed luctus arcu orci ac purus. Nulla tristique. Aliquam erat volutpat. Morbi in quam at neque egestas ultricies. Sed tempor consectetuer metus. Aliquam erat volutpat. Duis non tortor a quam tristique ultricies. Aenean eget augue eget nisi rutrum commodo. Vestibulum bibendum pellentesque lectus. Praesent et mauris a urna ultricies scelerisque. Class aptent taciti sociosqu ad litora torquent per conubia nostra, per inceptos himenaeos.

Nulla purus est, venenatis eget, hendrerit vitae, lacinia eget, dui. Maecenas tincidunt augue vel mi. Lorem ipsum dolor sit amet, consectetuer adipiscing elit. Fusce eu erat nec nisl pellentesque facilisis. Donec porttitor. Sed in nisl. Sed tristique tincidunt ipsum. Sed nec lectus. Quisque blandit, neque ut ullamcorper imperdiet, nibh justo imperdiet pede, sit amet egestas erat enim quis urna. In non urna eu lectus varius consequat. Nam interdum. Sed in metus vel enim ornare congue. Curabitur pede. Curabitur rutrum. Vestibulum ante ipsum primis in faucibus orci luctus et ultrices posuere cubilia Curae; Sed dui mauris, iaculis a, lobortis vitae, rhoncus id, nibh. Phasellus leo velit, dapibus quis, ultrices vitae, hendrerit sit amet, velit. Maecenas pretium pellentesque pede. Sed rutrum aliquet massa. Maecenas tellus.

\begin{sidewaystable}			
	\begin{center}
		\begin{tabular}{|c|c|c|c|c|c|c|c|c|c|c|c|c|c|}
			\hline
				\multirow{5}{1cm}{\scriptsize{\textbf{Tramo}}}	&
				\multirow{4}{0.15cm}{\scriptsize{\textbf{a}}}	&
				\multirow{4}{1.25cm}{\scriptsize{\textbf{Desnivel}}}	&
				\multicolumn{11}{|c|}{\scriptsize{Temperatura (\textdegree C)}}\\
			\cline{4-14}
				&
				&
				&
				\scriptsize{-5} &
				\scriptsize{5} &
				\scriptsize{10} &
				\scriptsize{15} &
				\scriptsize{20} &
				\scriptsize{25} &
				\scriptsize{30} &
				\scriptsize{35} &
				\scriptsize{40} &
					\scriptsize{45} &
					\scriptsize{50} \\
				\cline{4-14}
					&
					&
					&
					\multicolumn{11}{|c|}{\scriptsize{Tense (daN)}}\\
				\cline{4-14}
					&
					&
					&
					\scriptsize{$275.38$} &
					\scriptsize{$221.57$} &
					\scriptsize{$199.4$} &
					\scriptsize{$180.4$} &
					\scriptsize{$164.32$} &
					\scriptsize{$150.78$} &
					\scriptsize{$139.38$} &
					\scriptsize{$129.74$} &
					\scriptsize{$121.54$} &
					\scriptsize{$114.5$} &
					\scriptsize{$108.4$} \\
				\cline{4-14}
					&
					\scriptsize{(m)} &
					\scriptsize{(m)}	&
					\multicolumn{11}{|c|}{\scriptsize{Flecha (m)}}\\
				\hline
				\hline
					\scriptsize{1-2} &
					\scriptsize{80}	&
					\scriptsize{$11.09$}	&							
					\scriptsize{$0.55$} &
					\scriptsize{$0.68$} &
					\scriptsize{$0.75$} &
					\scriptsize{$0.83$} &
					\scriptsize{$0.91$} &
					\scriptsize{$1.00$} &
					\scriptsize{$1.08$} &
					\scriptsize{$1.16$} &
					\scriptsize{$1.24$} &
					\scriptsize{$1.31$} &
					\scriptsize{$1.39$} \\
						
				\hline								
		\end{tabular}
	\end{center}
	\caption[Ejemplo de tabla rotada]{Tabla rotada para aprovechar la altura de la p�gina.}
	\label{cuadro_regulacion_1_2}
\end{sidewaystable}

Donec a pede. Proin dolor. Ut nunc ligula, tempor id, ornare sit amet, aliquam et, nibh. In mollis iaculis pede. Vivamus gravida orci eu nisl. Sed nibh sem, consequat at, iaculis non, placerat in, ligula. Praesent id nisi. Nunc pellentesque justo non libero. Sed quis est sit amet purus lobortis blandit. Sed arcu justo, rhoncus condimentum, ullamcorper iaculis, viverra et, nisl. Mauris leo risus, placerat sit amet, commodo eget, lacinia at, ipsum. Fusce vulputate suscipit ipsum.



%
% SECCION - T�tulo de la secci�n
%
\section[�Nunc vel turpis ut varius vehicula est sed pulvinar?]{
	�Nunc vel turpis ut varius vehicula est?
}

Nunc vel turpis. Ut varius vehicula est. Sed pulvinar. Aliquam vehicula quam vitae sapien. Cras eget lectus a metus tempor venenatis. Phasellus et nulla. Class aptent taciti sociosqu ad litora torquent per conubia nostra, per inceptos himenaeos. Etiam sodales nibh aliquam dolor. Donec pulvinar lobortis risus. Ut augue erat, hendrerit a, accumsan non, elementum porttitor, est. Sed tincidunt neque nec ligula. Donec enim sapien, dictum gravida, porta non, molestie eu, lectus. Sed ultrices lectus at sem. Cras suscipit tortor id dolor. Nullam lacus risus, blandit non, lobortis sed, dignissim ut, nibh. Sed neque. Phasellus ullamcorper orci quis magna.

\begin{itemize}
	\item{
		�tem 1: Nam commodo ligula eget lorem. Nullam consectetuer nibh id orci. Aliquam ornare nisl nec risus. Quisque quis nunc. In ullamcorper, enim in eleifend aliquet, ipsum felis ultricies nibh, eu aliquam libero libero a lorem. Curabitur fringilla, massa nec vestibulum viverra, purus nulla ultricies ante, vitae dignissim tortor purus at dui. Fusce quis nisi. Aliquam erat volutpat. Sed molestie ultricies nunc. Nam vitae eros eget nibh tristique suscipit. Morbi ut velit. Quisque quis nulla a ipsum malesuada fermentum.
	}
	\item{
		�tem 2: Fusce scelerisque tortor ac libero. Class aptent taciti sociosqu ad litora torquent per conubia nostra, per inceptos himenaeos. Quisque nec est. Nunc ut risus a arcu congue ultricies. Aenean laoreet dapibus orci.
	}
	\item{
		�tem 3: Sed gravida dui et eros. Praesent id lacus. Etiam velit libero, tristique tempus, facilisis quis, bibendum euismod, magna. Quisque ut mi. Nam quis ipsum. Nulla porta dapibus arcu. Nullam ut purus vel libero scelerisque sodales.
	}
\end{itemize}

Aenean vitae sem nec turpis sagittis egestas. Sed nisi sapien, rutrum sit amet, dictum id, pretium eget, libero. Donec viverra risus eget mi. Duis porttitor. Nunc et sem eu augue dapibus condimentum. Mauris a tortor vel est adipiscing malesuada. Morbi neque mauris, vestibulum at, tincidunt vitae, convallis nec, ligula. Vestibulum quis turpis.



%
% SECCION - T�tulo de la secci�n
%
\section[Aenean sit amet lorem a dui sagittis tincidunt]{
	Aenean sit amet lorem a dui sagittis tincidunt
}
\label{ejemplo_etiqueta_para_referenciar_seccion_020}

Lorem ipsum dolor sit amet, consectetuer adipiscing elit. Nunc varius ornare mauris (ver cuadro \ref{ejemplo_referencia_a_tabla_020}). Donec orci quam, mattis eu, gravida in, condimentum eget, lacus. Integer a enim. Pellentesque turpis purus, feugiat in, dictum ut, interdum ut, nulla. Mauris purus.


% TODO: Tiene que haber alguna manera de hacer que la separaci�n entre la tabla y el p�rrafo sea configurado en el pre�mbulo...
\vspace{0.5cm}
\begin{table}[h]
	\begin{center}
		\begin{tabular}{lp{6cm}rr}
			\toprule
				% CABECERA
				Animal & 
				Description &
				Quantity & 
				Price (\$) \\
			\midrule
			\midrule
				% FILA 1
				Gnat &	
				Colloquial name to any of various small insects &
				10 & 
				13.65 \\
		
				% FILA 2 (con las dos primeras columnas en blanco)
				&
				&
				100 &
				6.00 \\
		
				% FILA 3
				Gnu &	
				An antelope of the genus Connochaetes &
				1 &
				92.50 \\
		
				% FILA 4
				Emu &
				The largest bird native to Australia &
				2 &
				33.33 \\
			
				% FILA 5
				Armadillo &
				Small placental mammals, known for having a bony armor shell &
				1 &
				8.99 \\
			\bottomrule
		\end{tabular}

	\end{center}

	\caption[Ejemplo de cuadro o tabla]{Ejemplo de cuadro o tabla.}
	\label{ejemplo_referencia_a_tabla_020}

\end{table}


% Subsecci�n
\subsection*{Donec consectetuer venenatis turpis}
Quisque enim neque, porta aliquam, tincidunt sed, eleifend quis, velit. Praesent sapien augue, suscipit nec, sodales a, faucibus in, mi. Ut et leo molestie massa elementum fringilla. Morbi sollicitudin sem vitae tortor tincidunt elementum. In ac nunc. Quisque placerat faucibus tellus. Nullam scelerisque, nibh et venenatis hendrerit, est diam semper mi, eu fringilla turpis augue vitae lectus. Ut accumsan mi at turpis. Sed sit amet magna eget augue vulputate blandit. Integer porttitor egestas nisl. Vestibulum consectetuer.

Quisque aliquet accumsan mauris. Nulla ante risus, semper eget, egestas sed, tristique vitae, arcu. Integer consequat dolor eget quam. Lorem ipsum dolor sit amet, consectetuer adipiscing elit. Donec consectetuer venenatis turpis. Etiam pharetra nisl vitae leo. Ut vehicula interdum leo. Donec in nisi. Morbi sagittis eleifend quam. Sed scelerisque mi at justo. Aliquam erat volutpat. Nulla consequat porttitor eros.


% Subsecci�n
\subsection*{Duis vestibulum risus vitae sem}
Suspendisse malesuada aliquam nunc. Integer neque. Nulla in magna. Duis vestibulum risus vitae sem. Morbi eu nunc sit amet erat facilisis pellentesque. Duis scelerisque risus vitae magna. Suspendisse placerat est eget lectus tristique tincidunt. Maecenas dolor nunc, vehicula in, vestibulum nec, varius a, sapien. Suspendisse in felis et ipsum tempor adipiscing. Cras elementum nulla vel lorem. Aliquam fringilla pharetra nulla. 

Aenean et lectus sed orci consequat faucibus. Nunc neque velit, feugiat at, iaculis at, ullamcorper nec, ligula. Sed laoreet, ipsum id iaculis bibendum, ipsum nibh bibendum urna, ut blandit arcu nisl blandit odio. Donec eros lorem, luctus at, ultrices sed, facilisis sed, est. Maecenas pharetra, felis a tempus aliquet, magna elit auctor lectus, eget euismod quam leo sit amet nibh. Ut commodo. Aliquam enim tellus, sodales vel, bibendum a, lacinia sed, metus. Nullam sed ligula ac libero rutrum molestie. Aliquam suscipit, eros non aliquet feugiat, magna mauris elementum pede, nec tristique nisl lorem hendrerit orci.


% Subsecci�n
\subsection*{Mauris molestie}
Aliquam vehicula purus id quam. Pellentesque egestas. Mauris molestie. Duis in arcu quis turpis elementum interdum. Curabitur eu orci. Nulla rhoncus convallis quam. Donec volutpat eleifend turpis. Ut nulla magna, hendrerit nec, laoreet ac, malesuada et, nisl. In dolor. Donec elit. Maecenas laoreet, eros non lobortis faucibus, turpis nisi viverra pede, vel ullamcorper pede tortor vel ligula.

Sed porttitor tincidunt mauris. Sed tempus libero sed turpis. Curabitur dictum adipiscing mi. Nunc porta leo quis nunc. Nunc euismod. Integer eu sem. Integer dapibus libero et eros. Mauris nisi. Quisque iaculis cursus enim (ver figura~\ref{ejemplo_referencia_a_figura_020}). Donec ut diam. Pellentesque erat.


% Incluimos una figura...
\begin{figure}[h]
	\centering
	\includegraphics[width=10cm]{./eps/control-pi.eps}
	\caption[Ejemplo de inclusi�n de figura o imagen]{
		Ejemplo de inclusi�n de figura o imagen.
	}
	\label{ejemplo_referencia_a_figura_020}
\end{figure}



%
% SECCION - T�tulo de la secci�n
%
\section[Praesent purus sem]{
	Praesent purus sem
}

Fusce fringilla feugiat turpis. Fusce bibendum. Quisque a magna. Aenean dignissim sollicitudin eros. Proin posuere vulputate nibh. Morbi nec diam vitae velit blandit tincidunt. Suspendisse in nisi sit amet ipsum luctus fringilla.


\paragraph*{Ut sed urna.}
Vivamus fermentum. Nulla eget lectus vel ligula dignissim condimentum. Fusce tristique, eros sit amet tempor congue, urna diam fermentum enim, ac ullamcorper urna magna ac ligula. Donec ultrices neque eu purus. Integer vehicula, mi eget mollis vehicula, risus tellus pulvinar dolor, a fringilla dui mauris sit amet leo. Nulla vestibulum est ut est. 

\paragraph*{Vivamus eu dui.}
Nulla tempus augue a sapien. Etiam non nibh. In lacus libero, ullamcorper eget, feugiat eu, tincidunt in, orci. Vivamus leo dolor, vehicula id, tempus et, venenatis at, dui. Pellentesque convallis ligula ut dui. Cras eu dolor sed dui tincidunt congue. Suspendisse justo elit, laoreet nec, fermentum id, lacinia viverra, lectus. Fusce erat velit, placerat id, dignissim et, gravida vitae, odio. Quisque massa enim, varius quis, pellentesque vel, hendrerit sit amet, mauris.

\paragraph*{etiam blandit imperdiet orci.}
Sed tempor pede a neque. Donec sed massa. Nam ornare auctor enim. Nunc vitae erat. In turpis. Quisque quam tortor, interdum sit amet, dapibus ac, bibendum et, sapien. In malesuada auctor elit. Pellentesque habitant morbi tristique senectus et netus et malesuada fames ac turpis egestas. Vivamus dignissim ligula quis ante. Phasellus dolor massa, accumsan in, aliquam nec, scelerisque ac, sem. Nunc ac orci. Donec risus libero, elementum non, euismod sollicitudin, pharetra ac, nibh. Curabitur nisi est, gravida vel, vulputate id, faucibus non, nibh.

\paragraph*{Curabitur iaculis ante.}
Phasellus nunc arcu, sodales eu, ultrices sed, suscipit ac, nisl. Quisque at velit. Suspendisse lobortis urna non metus. Donec lacinia congue risus. Suspendisse tempus nibh vel metus. Mauris lacinia volutpat nibh. Integer vitae mauris. Cras placerat lobortis nibh. In nisl odio, convallis a, vulputate id, ultrices in, nunc. 

\paragraph*{Morbi bibendum varius lorem.}
Curabitur eget sapien quis nulla fermentum iaculis. Vivamus elementum elit in ante. Nullam magna ipsum, semper quis, tempus sit amet, semper at, eros. In euismod molestie odio. Donec a ante. Morbi vel diam nec diam tincidunt luctus. Nunc rhoncus. Fusce metus justo, blandit facilisis, pulvinar eu, imperdiet sit amet, lectus.



%
% SECCION - T�tulo de la secci�n
%
\section[Vestibulum non lorem]{
	Vestibulum non lorem
}


Nam pellentesque, mauris eget mollis lobortis, lectus velit sollicitudin nulla, quis dapibus tortor mi in nibh. Cras viverra pede id augue. Ut sed urna ut pede rutrum vulputate. Sed pulvinar urna et eros. Proin cursus suscipit odio. Aenean aliquet. Nullam nibh ipsum, euismod facilisis, tincidunt eget, ultricies ut, turpis. Cras elementum. Proin vel eros vitae risus pretium aliquet. Sed sit amet purus in enim auctor iaculis. Curabitur pulvinar. Etiam in lectus sed enim egestas suscipit. Proin dignissim leo quis diam. Nam sed arcu eget velit vestibulum dapibus. 

\paragraph*{Integer ornare.}
Suspendisse lacinia sollicitudin mauris. Vestibulum felis est, sollicitudin sit amet, sollicitudin eget, dignissim eget, odio. Curabitur tincidunt metus euismod eros. Quisque libero.

\paragraph*{Duis semper dolor mattis quam.}
Praesent a ligula a enim sollicitudin laoreet. Cras varius semper eros. Aliquam egestas, sapien accumsan malesuada blandit, nisl magna viverra tortor, nec porttitor dolor nisi sed pede. Ut non nulla nec est aliquet porta. Nulla nec diam. Ut in libero vitae turpis fringilla posuere. Aliquam vestibulum pulvinar orci. Quisque lacinia neque porttitor nisi. Cum sociis natoque penatibus et magnis dis parturient montes, nascetur ridiculus mus. Sed augue diam, feugiat sed, pulvinar non, tempus vel, nunc. Nulla facilisi. Aenean sed enim sit amet odio aliquet dictum. Nulla vel arcu. Donec sodales. Nunc purus. Morbi semper lobortis nibh. Proin id risus ut elit hendrerit consectetuer. Vestibulum eu risus.

\paragraph*{Lorem ipsum}
Etiam ante risus, egestas ut, suscipit vel, auctor vitae, metus. Mauris fermentum laoreet est. Cum sociis natoque penatibus et magnis dis parturient montes, nascetur ridiculus mus. Aenean porttitor magna eget arcu aliquam dictum. 

Curabitur pulvinar. In tortor justo, malesuada at, ullamcorper a, tristique in, enim. Curabitur justo tellus, aliquet a, congue in, elementum quis, urna. Proin pulvinar enim a diam. Nunc adipiscing. Maecenas sodales elementum leo. Lorem ipsum dolor sit amet, consectetuer adipiscing elit. Proin et nisl non quam molestie consequat. 

Aliquam adipiscing bibendum justo. Duis mi. Maecenas semper eros at nibh. Maecenas pharetra odio et quam. Donec ligula diam, sodales condimentum, varius nec, lacinia sed, urna. Vestibulum dignissim diam quis lectus vehicula imperdiet. Etiam ornare, odio id lacinia semper, libero quam rhoncus eros, sollicitudin sagittis nulla lacus non elit. Etiam ultricies mi nec leo. 

\paragraph*{Curabitur eleifend libero sit amet odio.}
Curabitur sit amet orci at lacus euismod rutrum. Aliquam gravida venenatis nulla. Phasellus ligula. Praesent erat lorem, tincidunt id, ullamcorper et, egestas at, nulla. Cras sed dolor. Aliquam dignissim, nisl sit amet venenatis porttitor, erat nisl pharetra est, et hendrerit nisl odio nec elit. Cras consectetuer.


% subsecci�n
\subsection[Suspendisse sit amet elit]{Suspendisse sit amet elit}

Nam dui leo, suscipit ac, tincidunt ut, placerat eget, ipsum. Suspendisse feugiat volutpat ligula. Vestibulum ante ipsum primis in faucibus orci luctus et ultrices posuere cubilia Curae; Pellentesque vulputate molestie sem. Nullam porta odio et justo ultrices bibendum. Phasellus convallis vehicula mauris. Nunc cursus, elit in porta tincidunt, odio magna faucibus mi, ut pretium augue nisl a mi. Mauris quis massa. 

\paragraph*{Nunc pharetra.}
Nisl a convallis eleifend, pede urna aliquet lacus, ut condimentum tellus lectus eu arcu. Nam tincidunt luctus nulla. Pellentesque suscipit, justo sed interdum sodales, tellus elit dapibus tellus, sed ultrices magna arcu a magna. Duis velit diam, semper sit amet, venenatis sed, lacinia id, ligula. Integer malesuada. Cum sociis natoque penatibus et magnis dis parturient montes, nascetur ridiculus mus. Suspendisse et mi a sapien euismod sagittis. Nulla facilisi. Proin a felis imperdiet enim facilisis mollis. Nunc lacus lectus, scelerisque nec, pretium in, ullamcorper sed, quam.



%
% SECCION - T�tulo de la secci�n
%
\section[Integer ullamcorper nulla facilisi]{
	Integer ullamcorper nulla facilisi
}

Nulla porttitor, ipsum non gravida eleifend, dolor est accumsan nulla, sit amet ultrices libero felis vitae leo. Vestibulum at dui id metus porttitor pretium. Praesent et justo. Pellentesque habitant morbi tristique senectus et netus et malesuada fames ac turpis egestas. Morbi ipsum nisl, adipiscing et, tempus sit amet, luctus ornare, metus. Aliquam tellus sem, commodo blandit, vehicula et, vulputate congue, tellus. Cum sociis natoque penatibus et magnis dis parturient montes, nascetur ridiculus mus. Vivamus hendrerit. Proin tincidunt dui id felis.

Vestibulum eu justo a dolor aliquam ultrices. Aenean lobortis sollicitudin justo. Donec semper. Suspendisse eros erat, condimentum id, aliquam ac, hendrerit vel, eros. Proin ac nulla rutrum eros sagittis hendrerit. Phasellus aliquet, purus ut rhoncus bibendum, turpis libero vehicula ante, eu lacinia erat augue pulvinar lorem. Integer pede. Mauris congue dui vel nisl. Aenean sodales, nisl vitae fermentum sollicitudin, nisl massa mollis diam, sit amet ultricies velit quam vitae dui. Etiam accumsan eros ut tellus. Nullam arcu nunc, porttitor tincidunt, varius vitae, pellentesque nec, turpis. Morbi suscipit pellentesque sem. Donec faucibus. Praesent congue. Curabitur elit. Nullam sed enim nec sem sagittis dictum. Ut sed justo eu elit gravida tincidunt. Phasellus convallis, augue vel iaculis euismod, turpis nisl tincidunt tortor, ut scelerisque pede nisi id ligula. Suspendisse sagittis, orci ac eleifend gravida, orci augue sodales velit, sed iaculis felis dui sit amet nisi. Vivamus tincidunt euismod purus.

Aenean semper neque eu lectus. Donec mollis ante sit amet libero. Curabitur nec enim ut arcu suscipit cursus. Nulla ac erat et mauris sollicitudin volutpat. Curabitur ornare sapien ut purus. In faucibus sapien at neque. Suspendisse sed arcu et quam aliquet tristique. Praesent sit amet elit vitae turpis lobortis hendrerit. Aliquam nec mi iaculis mi tempor tincidunt. Donec at nisl convallis magna ullamcorper mollis. Etiam pretium erat eget metus adipiscing aliquam. Donec faucibus, sem viverra porta vestibulum, magna mi ultricies erat, aliquet lobortis dui sapien ut purus. Fusce tellus pede, sodales in, ultrices non, commodo non, lorem. Sed at magna scelerisque nisl sollicitudin imperdiet. Vestibulum commodo, nisi id bibendum consequat, velit tellus ultricies justo, at consectetuer eros nulla at massa. In a magna sit amet purus viverra aliquet. Morbi vel sem ut turpis commodo pretium. Pellentesque dolor sem, tincidunt id, sollicitudin et, tincidunt et, sem.

Duis quam libero, viverra sed, cursus nec, facilisis ut, nulla. Phasellus eros lectus, sollicitudin ac, egestas non, tristique eget, nibh. Sed pharetra, sem a mattis pretium, lectus lacus gravida nibh, sit amet porta tortor ipsum nec sapien. Aenean ipsum. Nulla sit amet justo. Aliquam aliquet nibh in erat. Quisque auctor est vitae nibh. Proin in nisi. Sed imperdiet. Etiam urna justo, lacinia ut, dapibus eu, commodo at, nibh. Pellentesque quis tellus at diam sagittis convallis. Integer ac sapien eget odio sagittis aliquet. Phasellus sagittis.

Etiam non leo. Aliquam quis dui. Suspendisse suscipit dui at lacus. Curabitur id urna sed purus ullamcorper ullamcorper. Nullam ut sem eu sapien malesuada rhoncus. In leo urna, rhoncus vel, tempus quis, ultricies vitae, ipsum. Sed est justo, sodales a, semper et, pretium in, massa. Etiam aliquam est a ligula. Nullam eleifend scelerisque nisi. In sollicitudin mi vel nibh. Praesent dapibus, nisl quis tincidunt dictum, lorem quam dapibus lorem, tristique mollis mi tortor ut nulla. Sed sit amet sapien. Aenean non urna lacinia urna euismod interdum. Praesent in nibh. In mauris. Ut scelerisque lectus vitae neque.



%
% FIN DEL CAP�TULO
%
	%\include{pfc_main_030}
	%\include{pfc_main_040}
	%\include{pfc_main_050}
	%\include{pfc_main_060}


	\backmatter
	% INCLUIMOS LOS APÉNDICES... !! DELETED - INCLUDE IF IT'S NECESSARY
        %\appendix
	%%
% Memoria económica.
%

\chapter[Título de apéndice]{
	Título de apéndice.
}

Lorem ipsum dolor sit amet, consectetur adipiscing elit. Integer luctus varius augue vel imperdiet. Maecenas neque turpis, accumsan ut eleifend vel, eleifend gravida est. Duis mattis pulvinar porta. Nunc consequat, felis ut malesuada egestas, ipsum urna sodales enim, eu lobortis mauris tellus in augue. In tristique tortor id nisi placerat et rutrum tortor auctor. Aenean erat elit, faucibus non egestas semper, volutpat sit amet ante. Nam nulla justo, suscipit nec sagittis sed, varius id magna. Duis a lorem id felis eleifend aliquet eget vel arcu. Etiam ultricies diam a magna facilisis ornare. Integer tempor massa quis sem gravida bibendum. Sed massa felis, condimentum sit amet bibendum in, tristique sit amet mauris. Nullam tellus ante, ornare eget mollis ac, gravida id nunc. Donec elit neque, aliquam eget auctor sed, gravida in nibh. Vestibulum sit amet diam a magna tempor bibendum quis a ante. Integer augue sapien, facilisis ac malesuada sed, consectetur id turpis. Quisque a leo dapibus felis aliquet vestibulum vitae vitae urna. Nunc ut felis ut metus gravida ultricies. Cras tincidunt nisi ante. Nunc eu quam interdum dui cursus lacinia. Praesent adipiscing ultrices nulla at posuere.

Aliquam gravida facilisis eleifend. Nunc id lorem nibh. Nulla facilisi. Nunc suscipit lectus quis tellus adipiscing adipiscing. Proin ipsum justo, egestas dictum iaculis et, pretium non diam. Aenean vehicula hendrerit lectus, vel bibendum nisi volutpat quis. Fusce molestie condimentum libero, nec scelerisque nisi mollis et. Aliquam placerat accumsan congue. In scelerisque nulla sit amet mauris consectetur consectetur. Suspendisse mattis mi in sapien egestas sit amet accumsan urna sagittis. Mauris metus massa, porta at rutrum eget, sollicitudin eu risus. Donec non diam erat. Nam sagittis, nibh non fermentum ornare, ipsum est mollis felis, eget vulputate tellus ante at erat. Praesent porta tempus velit at faucibus. Donec rhoncus varius nunc, eget imperdiet ligula blandit eget. Morbi pulvinar, nunc ac porta interdum, mauris ligula sollicitudin sem, et sollicitudin turpis magna eu lectus. Pellentesque nulla dui, feugiat quis pharetra eu, accumsan imperdiet quam. Aenean lacinia porttitor leo vitae hendrerit. Pellentesque dictum ultricies dui, et rhoncus nulla sagittis ut. Suspendisse suscipit, risus sed posuere rhoncus, neque metus scelerisque diam, eu gravida quam orci eu dui.

Nulla sit amet sem massa, aliquet auctor magna. Vivamus quis magna sollicitudin felis bibendum mollis. Etiam elementum lacus id quam tincidunt pulvinar. Ut vehicula aliquam ante, vel venenatis libero bibendum eget. Sed gravida lobortis velit, et ornare nunc ullamcorper et. Morbi neque sapien, mattis id semper in, varius at arcu. Duis aliquam ornare sem, eu sagittis libero facilisis non. Nunc quis ultricies augue. Pellentesque ac nulla libero, ut suscipit dui. Morbi aliquet quam non metus iaculis quis accumsan ipsum consectetur.

Mauris mauris odio, varius quis accumsan bibendum, vulputate vel tellus. Sed non consectetur eros. Aenean et mi ut est tincidunt commodo. Donec nunc metus, facilisis a consequat vitae, dignissim quis leo. Cras condimentum, felis a molestie viverra, massa eros consequat nisl, at porta nibh sapien condimentum ante. Duis convallis velit quis metus pulvinar convallis. Suspendisse sed purus non purus dictum ultricies in vitae nisl. Cras convallis erat eu ipsum molestie fermentum. Mauris turpis nunc, condimentum non aliquam ac, ornare at lectus. Donec scelerisque fermentum malesuada. Praesent urna nisl, iaculis sed pharetra eu, tincidunt et lectus. Nulla posuere eros at elit vehicula laoreet. Sed ut erat orci, ac pellentesque nibh. Nulla at sapien sed justo luctus laoreet.

Sed convallis massa laoreet mi venenatis eget dapibus orci varius. Pellentesque purus nisi, placerat nec pretium ac, pretium vitae augue. Vivamus eleifend vestibulum mi, ac dignissim urna vulputate vel. Integer volutpat neque purus. Nulla hendrerit tempus est nec aliquet. Ut facilisis bibendum nisl eget porttitor. Cum sociis natoque penatibus et magnis dis parturient montes, nascetur ridiculus mus. Fusce consectetur sem eget ante rhoncus nec mollis lorem volutpat. Fusce in libero lectus, ut gravida lectus. Nam nulla nunc, dignissim sed sollicitudin a, commodo vel felis. Donec lacinia justo quis neque tempor interdum ut cursus mi. Phasellus venenatis porttitor mi consectetur egestas. Pellentesque quis elit nibh. Pellentesque molestie turpis ac arcu gravida sed ultricies felis pulvinar.

%         \include{pfc_appendix_020}

	
	% INCLUIMOS LA BIBLIOGRAFÍA...
	\nocite{*}	% Se usa para indicar en la bibliografía las referencias no citadas.
	\bibliography{biblio}
	\bibliographystyle{plain}

\end{document}

